% Preamble of theme-report/main.tex, reused unchanged (blank-line runs collapsed).
\documentclass[12pt, a4paper, oneside, final]{book} 
\usepackage{cmap}
\usepackage[unicode]{hyperref} 
\usepackage[utf8]{vietnam}
\usepackage{amsthm}
\usepackage{yhmath}
\usepackage{amsmath}
\usepackage{mathpazo}
\usepackage{booktabs}
% \usepackage{xcolor}

\usepackage{listings}
\usepackage{xcolor}

% Định nghĩa màu nền cho ô code
\definecolor{codegray}{rgb}{0.95,0.95,0.95}

\lstdefinestyle{pythonstyle}{
    language=Python,
    backgroundcolor=\color{codegray},
    basicstyle=\ttfamily\small,
    keywordstyle=\color{blue},
    stringstyle=\color{red},
    commentstyle=\color{green!60!black},
    showstringspaces=false,
    breaklines=true,
    frame=single,
    columns=flexible
}

\lstdefinestyle{sqlstyle}{
    language=SQL,
    basicstyle=\ttfamily\small,
    keywordstyle=\color{blue}\bfseries,
    commentstyle=\color{gray},
    stringstyle=\color{red!70!black},
    showstringspaces=false,
    breaklines=true,
    frame=single,
    columns=flexible
}

% \usepackage{arev}
% \usepackage[charter]{mathdesign}
% \usepackage[utopia]{mathdesign}
% \usepackage[mathscr]{eucal}
% \usepackage{fourier}
% \usepackage{charter}
% \usepackage{ccfonts} 
\usepackage{tasks}
\usepackage[dvipsnames]{xcolor}
\usepackage{tikz}
\usetikzlibrary{calc,angles}
\usetikzlibrary{intersections}
\usepackage{pgfplots}
\usepackage{indentfirst}
\usepackage{enumitem}
\usepackage{commath}
\usepackage{scrextend}
%%%%%%%%%%%%%%%%%%%%%%%%%%%%%%%%%%%%%%%%%%%
\usepackage{tikz,tkz-tab}
\usetikzlibrary{calc, angles, quotes, intersections, positioning, patterns}
\usepackage{pgfplots}
%%%%%%%%%%%%%%%%%%%%%%%%%%%%%%%%%%%%%%%%%%%
\usepackage{nicefrac}
\usepackage{picinpar}
% \usepackage{slashbox}
\usepackage{amssymb}
\usepackage{multicol} 
\usepackage{fancyhdr} 
\usepackage{graphicx}
\usepackage{geometry}
\usepackage{wrapfig}
\usepackage{multirow}
\usepackage{tabularx}
%\usepackage{standalone}%chen hinh bang file tex
%\usepackage{pstricks-add}
%\usepackage{pstricks,pst-plot,pst-text,pst-tree,pst-eps,pst-fill,pst-node,pst-math}

%\usepackage{mathptmx}  %Bổ sung thêm dòng này để dùng watermark
%\usepackage{draftwatermark} %Bổ sung thêm dòng này để dùng watermark

\geometry{left=2.5cm,right=1.5cm,top=2cm,bottom=2cm}

\setlength{\oddsidemargin}{-0.1in}
\setlength{\evensidemargin}{0.0in}
\setlength{\textwidth}{6.8in}
\setlength{\topmargin}{-0.5in}
\setlength{\textheight}{10.0in}
%%%%%%%%%%Mau moi
\colorlet{NAVYBLUE}{NavyBlue}
\usepackage{sectsty}
%\renewcommand{\thesection}{}

\sectionfont{\centering \fontsize{16}{16} \color{orange}\selectfont \scshape \MakeUppercase}

% Define rule (color) for header and footer
%\renewcommand{\headrule}{{\color{GreenYellow}\rule{\textwidth}{0.1pt}}} 
%\renewcommand{\footrule}{{\color{GreenYellow}\rule{\textwidth}{0.1pt}}} 
%%%%%%%%%%%%%%

\pagestyle{fancy} 
\fancyhead[L]{\texttt{}}
\fancyhead[R]{\texttt{\textcolor{NavyBlue} {\it }}} 
\fancyfoot[C]{} 
\fancyfoot[L]{ \textsf{\rightmark}}
\fancyfoot[R]{\textsf{\thepage}}
\renewcommand{\chaptermark}[1]{%
	\markboth{\chaptername\ \thechapter.\ #1}{}} %bỏ đi sẽ thấy chữ in
\renewcommand{\sectionmark}[1]{%
	\markright{\thesection.\ #1}} %bỏ đi sẽ thấy chữ in

\theoremstyle{definition} 
\newtheorem*{proposition}{\color{NavyBlue} Proposition}
\newtheorem*{definition}{\color{NavyBlue} Definition} 
\newtheorem{exercise}{\color{NavyBlue} Bài tập}[chapter] 
\newtheorem{bt}{\color{NavyBlue} Bài}[chapter]
%\newtheorem{exercise}{\color{blue} Bài tập}[chapter]
\newtheorem{tb}{\color{NavyBlue} Bảng}%[chapter]
\newtheorem{BT}{\color{NavyBlue} Bài toán}[chapter]
\newtheorem{vd}{\color{NavyBlue} Ví dụ}[chapter]
\newtheorem{dl}{\color{NavyBlue} Định lý}[chapter]
\newtheorem{nl}{\color{NavyBlue} Nguyên lý}[chapter]
\newtheorem{hq}{\color{NavyBlue} Hệ quả}[chapter]
\newtheorem{dn}{\color{NavyBlue} Định nghĩa}[chapter]
\newtheorem{tc}{\color{NavyBlue} Tính chất}[chapter]
\newtheorem{kn}{\color{NavyBlue} Khái niệm}[chapter]
\newtheorem{nx}{\color{NavyBlue} Nhận xét}[chapter]
\newtheorem{qt}{\color{NavyBlue} Quy tắc}[chapter]
\newtheorem{gt}{\color{NavyBlue} Giả thiết}[chapter]
\newtheorem{pl}{\color{NavyBlue} Phụ lục} %[chapter]
\newtheorem{cy}{\color{NavyBlue} Chú ý}[chapter]
\newcommand{\Lga}{\textcolor{NavyBlue}{\it Lời giải. \ }}
\newcommand{\Sol}{\textcolor{NavyBlue}{\it Lời giải. \  }}
\newcommand{\Gy}{\textcolor{NavyBlue}{\it Gợi ý. \  }}
\newcommand{\chm}{{\it Chứng minh. }}
\newcommand{\hd}{\textcolor{NavyBlue}{\it Hướng dẫn. }}
\newcommand{\ds}{\textcolor{NavyBlue}{\it \noindent Đáp số. }}
\newtheorem{example}{Ví dụ}[section]
\newcommand{\eproof}{\hfill $\square$}
\theoremstyle{plain} 
\newtheorem*{theorem}{Định lý}

\theoremstyle{remark} 
\newtheorem*{remark}{Chú ý}

\providecommand{\lr}[1]{\left( #1 \right)} 
\providecommand{\Lr}[1]{\left[ #1 \right]}
\providecommand{\LR}[1]{\left\{ #1 \right\}} 
\providecommand{\abs}[1]{\left\vert #1 \right\vert} 
\providecommand{\norm}[1]{\left\Vert #1 \right\Vert} 
\providecommand{\e}{\varepsilon} 
\providecommand{\Sum}[2]{\sum\limits_{#1}^{#2}} 
\providecommand{\Lim}{\lim\limits_{n \to +\infty}} 
\providecommand{\SUM}{\Sum{n = 1}{\infty}} 
\providecommand{\X}{\mathbb{X}}
\providecommand{\N}{\mathbb{N}}
\providecommand{\Z}{\mathbb{Z}}
\providecommand{\Q}{\mathbb{Q}}
\providecommand{\I}{\mathbb{I}}
\providecommand{\R}{\mathbb{R}}
\providecommand{\D}{\mathbb{D}}
\providecommand{\e}{\varepsilon}
\providecommand{\d}{\partial}
\providecommand{\ov}[1]{\overline{#1}} 
\providecommand{\al}[1]{\alpha{#1}} 
\providecommand{\ud}[1]{\underline{#1}} 

\DeclareMathSymbol{\INT}{\mathop}{largesymbols}{"5A} 
\renewcommand{\int}{\INT} 
\renewcommand{\iint}{\mathop{\int\!\!\!\!\!\int}}
\renewcommand{\iiint}{\mathop{\int\!\!\!\!\!\int\!\!\!\!\!\int}}
\DeclareMathOperator{\Int}{int} 
\DeclareMathOperator{\grad}{grad} 
\DeclareMathOperator{\sgn}{sgn} 
\providecommand{\hcn}{hình chữ nhật }
\DeclareMathOperator{\rank}{rank}
\DeclareMathOperator{\arccot}{arccot}
\DeclareMathOperator{\const}{const}
\DeclareMathOperator{\tg}{tg}
\DeclareMathOperator{\arctg}{arctg}
\DeclareMathOperator{\cotg}{cotg}
\DeclareMathOperator{\arccotg}{arccotg}
\DeclareMathOperator{\sh}{sh}
\DeclareMathOperator{\ch}{ch}
\DeclareMathOperator{\Sup}{Sup}
\DeclareMathOperator{\Inf}{Inf}
\DeclareMathOperator{\Min}{Min}
\DeclareMathOperator{\Max}{Max}

\newcounter{exerenumi}
\renewcommand{\theexerenumi}{\bf\alph{exerenumi}}
\newcommand{\labelexerenumi}{\theexerenumi.}
\newenvironment{answer}{\begin{list}{\labelexerenumi}{%
			\usecounter{exerenumi}%
			\setlength{\labelwidth}{3pt}%
			\setlength{\leftmargin}{0pt}%
			\setlength{\itemindent}{15pt}}}{\end{list}}
\def\baselinestretch{1.3}

\usepackage{fancyhdr}

\usepackage{xcolor}
\usepackage{sectsty}
\chapterfont{\color{NavyBlue}}  % sets colour of chapters
\sectionfont{\color{NavyBlue}}  % sets colour of sections
\subsectionfont{\color{NavyBlue}}  % sets colour of sections


